\documentclass[12pt]{article}

\usepackage{amsmath,amsthm,amssymb}
\usepackage[portuguese]{babel}
\usepackage{indentfirst}
\usepackage{titlesec}

\titleformat*{\section}{%
    \fontsize{14}{16}\bfseries%
}

\titleformat*{\subsection}{%
    \fontsize{10}{12}\bfseries%
}

\title{Relatório - Lista Computacional}
\author{Gabriel Falcão, Marcelo Lopes, Mateus Almeida, Vinícius Kern\\ 
MC22 - Otimização Linear (2026.2)}
\date{20 de setembro de 2026}

\begin{document}

\maketitle

\section*{Parte 1: Auditoria e Depuração de Código}
Ao final da primeira fase do algoritmo, é esperado que seja encontrada uma base ótima para o problema, eliminando as variáveis artificiais introduzidas previamente. 

Com o valor ótimo sendo obtido durante a primeira fase, é garantido que o problema é viável. Porém, se alguma variável artificial permanecer na base com valor nulo, existem duas possibilidades:

1. A matriz A possui restrições redundantes;

2. O problema é degenerado.

No primeiro caso, é necessário descartar as colunas artificiais enquanto a matriz de base $B$ é mantida quadrada e com vetores linearmente independentes, pois deve ser não-singular para a aplicação do algoritmo na segunda fase, que exige multiplicações por $B^{-1}$.

A fim de lidar de forma adequada com as colunas artificiais restantes, foi implementado o método de pivoteamento neutro.

\section*{Parte 2: Instância Parametrizada por Matrícula}
O problema de programação linear a ser montado, em função da matrícula de um dos estudantes, é dado por:
$$
\min_{x\geq0}-(d_1+1)x_1-(d_2+2)x_2-(d_3+1)x_3
$$
$$\text{s.a.}\quad
\begin{pmatrix}
2 & 1 & 1\\
1 & 2 & d_4+1\\
1 & 1 & 2
\end{pmatrix}
\begin{pmatrix}
x_1\\x_2\\x_3
\end{pmatrix}
\leq
\begin{pmatrix}
10+d_1\\12+d_2\\8+d_3
\end{pmatrix}
$$

Utilizando os dígitos da matrícula $d_1=3,d_2=3,d_3=5,d_4=0$, obtivemos o seguinte problema linear:
$$
\min_{x\geq0}-4x_1-5x_2-6x_3
$$
$$\text{s.a.}\quad
\begin{pmatrix}
2 & 1 & 1\\
1 & 2 & 1\\
1 & 1 & 2
\end{pmatrix}
\begin{pmatrix}
x_1\\x_2\\x_3
\end{pmatrix}
\leq
\begin{pmatrix}
13\\15\\13
\end{pmatrix};\,
\begin{pmatrix}
x_1\\x_2\\x_3
\end{pmatrix}
\geq0
$$

Por meio do código corrigido, foi possível encontrar a solução ótima, com as variáveis originais ($x_1,x_2,x_3$) na base, como esperado, e obtivemos os seguintes resultados computacionalmente:

Vetor de custos reduzidos final:
$$\overline{c}=c - A^Ty=
\begin{pmatrix}
        0 & 0 & 0 & 0.2500 & 1.2500 & 2.2500
\end{pmatrix}^T
$$

Inversa da base ótima:
$$B^{-1}=
\begin{pmatrix}
0.7500 & -0.2500 & -0.2500\\
-0.2500 & 0.7500 & -0.2500\\
-0.2500 & -0.2500 & 0.7500
\end{pmatrix}
$$

Vetor dual:
$$y^*=(B^{-1})^Tc_B=
\begin{pmatrix}
-0.25 \\ -1.25 \\ -2.25
\end{pmatrix}
$$

Para o intervalo de variação do coeficiente $c_1$, realizamos uma análise manual. O custo reduzido ao realizar uma perturbação $\delta$ no valor inicial $c_1=-4$ deve manter o custo reduzido das variáveis fora da base maior que ou igual a zero, ou seja:

$$\overline{c}_N=c_N-N^Ty_{novo}$$
$$y_{novo}=B^{-1}\left(c_B+\begin{pmatrix}\delta\\0\\0\end{pmatrix}\right)=y^*+B^{-1}\begin{pmatrix}\delta\\0\\0\end{pmatrix}$$

Note que $N=I$ e $c_N=0$, pois as variáveis não-básicas são exatamente as variáveis de folga introduzidas. Assim, realizando os cálculos em função de $\delta$, temos:

$$y_{novo}=
\begin{pmatrix}
-0.25\\-1.25\\-2.25
\end{pmatrix}
+
\begin{pmatrix}
0.75 & -0.25 & -0.25\\
-0.25 & 0.75 & -0.25\\
-0.25 & -0.25 & 0.75
\end{pmatrix}
\begin{pmatrix}
\delta\\0\\0
\end{pmatrix}
=
\begin{pmatrix}
-0.25+0.75\delta\\-1.25-0.25\delta\\-2.25-0.25\delta
\end{pmatrix}
$$

$$\overline{c}_N=-
\begin{pmatrix}
-0.25+0.75\delta\\-1.25-0.25\delta\\-2.25-0.25\delta
\end{pmatrix}
=
\begin{pmatrix}
0.25-0.75\delta\\1.25+0.25\delta\\2.25+0.25\delta
\end{pmatrix}
$$

Assim, temos as seguintes restrições:

$$0.25-0.75\delta\geq0\Rightarrow\delta\leq\frac{1}{3}$$
$$1.25+0.25\delta\geq0\Rightarrow\delta\geq-5$$
$$2.25+0.25\delta\geq0\Rightarrow\delta\geq-9$$

Assim, $\delta\in\left[-5,\frac{1}{3}\right]$

Portanto, obtemos o seguinte intervalo para $c_1$, com base no valor original $-4$ e nos limites para $\delta$:
$$c_1\in\left[-9,\frac{11}{3}\right]$$

\end{document}